\documentclass{article}
\usepackage[utf8]{inputenc}
\usepackage[T1]{fontenc}
\usepackage{amsmath,amssymb,amsfonts}
\usepackage{booktabs}
\usepackage{multirow}
\usepackage{graphicx}
\usepackage{hyperref}
\usepackage[margin=1in]{geometry}
\usepackage{xcolor}
\usepackage{natbib}
\usepackage{caption}
\usepackage{subcaption}

\title{Cross-Lingual Conflict Detection in Multilingual LLMs:\\
Mechanistic Evidence for English-Anchored Representations}

\author{
Anonymous Authors
}

\date{}

\begin{document}

\maketitle

% ============================================================
% ABSTRACT
% ============================================================
\begin{abstract}
Multilingual large language models resolve entity-centric knowledge conflicts in ways that depend on linguistic affinity rather than resource abundance, yet the internal mechanisms underlying this behavior remain unexplored. We present a mechanistic investigation of cross-lingual conflict detection signals in the residual stream of three transformer architectures (Llama-3.1-8B, Qwen3-8B, Gemma-2-9B-IT) across two typologically distinct language pairs (EN$\leftrightarrow$ZH, EN$\leftrightarrow$JA). Linear probes trained on residual stream activations detect conflict signals at 70--77\% balanced accuracy, with cross-lingual transfer exhibiting a consistent directional asymmetry: transfer toward English (X$\to$EN) outperforms transfer away from English (EN$\to$X) by 6.3--7.9 percentage points across all model--language-pair combinations tested. Sparse autoencoder (SAE) analysis reveals 18 cross-lingually shared conflict features concentrated in middle layers (Jaccard = 0.28) that diverge in late layers (Jaccard = 0.12, $p = 0.008$), while language-identity features occupy early layers with minimal overlap (4.3\%). Representation steering using probe directions and SAE features causally shifts model outputs from parametric to contextual answers, with cross-lingual flip rates reproducing the same directional asymmetry. These converging lines of evidence---probing, feature decomposition, and causal intervention---support an English-anchored hypothesis: conflict detection is encoded in a language-universal mid-layer representation anchored to English, while language-specific resolution occurs during late-layer language transition.
\end{abstract}

% ============================================================
% 1. INTRODUCTION
% ============================================================
\section{Introduction}

When multilingual LLMs encounter factual knowledge conflicts---cases where contextual information contradicts parametric knowledge---the language of presentation systematically influences which knowledge source prevails. Recent work by Zhao et al.\ (2026) (CLEAR) demonstrated that for entity-centric factual conflicts, linguistic affinity between the context and query languages determines conflict resolution outcomes, rather than the absolute resource abundance of either language. This finding raises a fundamental mechanistic question: \emph{how are cross-lingual conflict signals encoded in the model's internal representations, and why does linguistic similarity play a decisive role?}

We investigate this question through three complementary analysis tracks applied to the residual stream of multilingual transformers. \textbf{Track 1} (probing) trains linear and MLP classifiers on residual stream activations to detect and transfer conflict signals across languages. \textbf{Track 2} (SAE decomposition) uses pre-trained sparse autoencoders to identify monosemantic conflict features and characterize their cross-lingual sharing patterns. \textbf{Track 3} (causal intervention) steers model behavior by manipulating conflict-related directions in activation space.

We formulate two competing hypotheses about the encoding of cross-lingual conflict signals:

\begin{itemize}
    \item \textbf{H1 (English-anchored universal detection)}: Conflict detection signals are language-independent, encoded in a mid-layer shared representation space anchored to English. Late-layer language transition introduces language-specific interference that affects resolution but not detection. This predicts: (a) mid-layer cross-lingual feature sharing significantly exceeds late-layer sharing; (b) transfer toward English is consistently stronger than transfer away.

    \item \textbf{H2 (Language-pair-specific encoding)}: Conflict features are language-pair-specific from mid-layers onward, with linguistic distance directly encoded in conflict representations. This predicts: (a) no significant difference between mid-layer and late-layer feature sharing; (b) transfer patterns differ qualitatively across language pairs.
\end{itemize}

Our experiments on three architectures (Llama-3.1-8B, Qwen3-8B, Gemma-2-9B-IT) and two language pairs (EN$\leftrightarrow$ZH, EN$\leftrightarrow$JA) yield four principal findings:

\begin{enumerate}
    \item Cross-lingual conflict signals are linearly detectable in the residual stream at 64--77\% balanced accuracy across all directions and models, with permutation test $p < 0.001$.
    \item Cross-lingual transfer is directionally asymmetric: X$\to$EN exceeds EN$\to$X by 6.3--7.9pp across all 4 model$\times$language-pair combinations (3/4 statistically significant at 95\% CI, 4/4 directionally consistent).
    \item SAE conflict features are cross-lingually shared in middle layers (Jaccard = 0.28) but diverge in late layers (Jaccard = 0.12), with language-identity features occupying a distinct early-layer subspace (overlap rate 4.3\%).
    \item Representation steering causally shifts conflict resolution behavior, with cross-lingual flip rate asymmetry mirroring the probe transfer asymmetry.
\end{enumerate}

These results collectively support H1: conflict detection is encoded in an English-anchored language-universal representation, with language-specific processing emerging during late-layer transition.

% ============================================================
% 2. RELATED WORK
% ============================================================
\section{Related Work}

\paragraph{Knowledge conflicts in LLMs.}
Knowledge conflicts arise when contextual information contradicts a model's parametric knowledge. Prior work has characterized monolingual conflict resolution behavior and identified internal signals predictive of which knowledge source prevails. However, the cross-lingual dimension remains largely unexplored. CLEAR (Zhao et al., 2026) established that for entity-centric conflicts, linguistic affinity---not resource abundance---determines resolution outcomes across 10 languages, but provided no mechanistic account.

\paragraph{Cross-lingual representations in multilingual models.}
Multilingual transformers develop shared cross-lingual representations in middle layers before diverging in late layers for language-specific decoding. Wendler et al.\ (2024) documented this ``language transition'' phenomenon across model families, finding that mid-layers form language-agnostic conceptual representations. Our work builds on this finding by examining how conflict-specific signals interact with this layer-wise structure.

\paragraph{Sparse autoencoders for mechanistic interpretability.}
Sparse autoencoders (SAEs) decompose neural activations into interpretable monosemantic features. Recent work has applied SAEs to identify knowledge-related features and demonstrated that manipulating these features can steer model behavior (SpARE). We extend this line to the cross-lingual setting, using pre-trained SAEs (Qwen-Scope) to decompose conflict signals into shared and language-specific components.

\paragraph{Representation steering.}
Activation steering modifies model behavior by adding direction vectors to the residual stream. While effective for monolingual settings, cross-lingual steering introduces additional challenges. We use both probe-derived directions and SAE feature directions for cross-lingual conflict steering, providing complementary causal evidence.

% ============================================================
% 3. METHOD
% ============================================================
\section{Method}

\subsection{Data Construction}

We construct cross-lingual knowledge conflict datasets from ConflictQA, which contains entity-centric factual questions where contextual and parametric answers diverge. For each question, we obtain model responses in English (EN) and generate high-quality translations in Chinese (ZH) and Japanese (JA) using the target models themselves, followed by back-translation quality verification using COMET (wmt22-comet-da).

\paragraph{Translation quality validation.}
Back-translation COMET scores are computed for all translated samples. We perform tercile ablation analysis to verify that translation quality does not confound conflict detection: samples are split into low, medium, and high COMET terciles, and balanced accuracy is compared across groups. As shown in Table~\ref{tab:comet}, both Llama-3.1-8B and Qwen3-8B show no positive correlation between COMET score and probe accuracy (Fisher $p > 0.69$), with the high-COMET tercile actually producing \emph{lower} accuracy than the low-COMET tercile. This rules out translation quality as a confound.

\begin{table}[t]
\centering
\caption{COMET tercile ablation. Probe accuracy does not increase with translation quality, ruling out translation quality as a confound.}
\label{tab:comet}
\small
\begin{tabular}{llcccc}
\toprule
Model & Tercile & $N$ & COMET$_{\text{mean}}$ & Bal.\ Acc. & $\Delta$ vs Low \\
\midrule
\multirow{3}{*}{Llama-3.1-8B} & Low & 91 & 0.582 & 64.8\% & --- \\
 & Med & 91 & 0.714 & 66.1\% & +1.3pp \\
 & High & 90 & 0.821 & 63.2\% & $-$1.6pp \\
\midrule
\multirow{3}{*}{Qwen3-8B} & Low & 56 & 0.585 & 66.4\% & --- \\
 & Med & 56 & 0.700 & 67.7\% & +1.3pp \\
 & High & 56 & 0.816 & 60.7\% & $-$5.6pp \\
\bottomrule
\end{tabular}
\end{table}

\subsection{Track 1: Residual Stream Probing}

For each model, we extract residual stream activations at every layer for conflict and non-conflict samples in each language. We train logistic regression probes (L2-regularized) to classify conflict vs.\ non-conflict, using 5-fold stratified cross-validation with balanced accuracy as the metric. Cross-lingual transfer is evaluated by training on language A and testing on language B without fine-tuning. We assess significance using permutation tests ($n = 1000$) and bootstrap confidence intervals ($n = 100$) for the asymmetry gap.

MLP probes (single hidden layer, 256 units) are trained with identical splits to quantify non-linear components. Statistical significance is determined at $\alpha = 0.05$.

\subsection{Track 2: SAE Feature Analysis}

We apply pre-trained sparse autoencoders from Qwen-Scope to Qwen3-8B-Base activations. Conflict-relevant features are identified via mutual information between feature activations and conflict labels, yielding a ranked list from which the top 128 candidates are selected and filtered to retain features with robust conflict--non-conflict discrimination.

To disentangle conflict features from language-identity features, we apply LangFIR-style random-token filtering: features that activate consistently across random token sequences in a given language are classified as language-identity features, and their overlap with conflict features is quantified.

Cross-lingual feature sharing is measured by the Jaccard index of EN and ZH conflict feature sets, computed separately for early (0--33\%), mid (33--67\%), and late (67--100\%) layer groups. The H1 vs.\ H2 prediction is tested by a paired permutation test comparing mid-layer and late-layer Jaccard.

\subsection{Track 3: Causal Intervention via Steering}

\paragraph{Probe-based steering.}
The weight vector of the best-layer linear probe defines a conflict direction in activation space. We add $\alpha \cdot \mathbf{v}$ to the residual stream at the optimal layer (where $\alpha$ is tuned to maximize flip rate while preserving coherence) and measure the rate at which model outputs shift from parametric to contextual answers.

\paragraph{SAE feature steering.}
We construct steering vectors from SAE conflict features (shared, EN-only, ZH-only, and combined sets) by clamping selected feature activations. Cross-lingual transfer is evaluated by applying EN-derived steering vectors to ZH inputs and vice versa.

% ============================================================
% 4. EXPERIMENTS AND RESULTS
% ============================================================
\section{Experiments and Results}

\subsection{Cross-Architecture Conflict Detectability}
\label{sec:detectability}

We train linear probes on residual stream activations from three model architectures: Llama-3.1-8B, Qwen3-8B, and Gemma-2-9B-IT. Table~\ref{tab:cross_arch} reports balanced accuracy for all four transfer directions (within-language and cross-lingual) along with the asymmetry gap and bootstrap confidence intervals.

\begin{table}[t]
\centering
\caption{Cross-architecture probe results. EN$\to$EN and ZH$\to$ZH are within-language; EN$\to$ZH and ZH$\to$EN are cross-lingual transfer. Asym.\ gap = (ZH$\to$EN) $-$ (EN$\to$ZH). All permutation test $p < 0.001$.}
\label{tab:cross_arch}
\small
\begin{tabular}{lccccccl}
\toprule
Model & EN$\to$EN & ZH$\to$ZH & EN$\to$ZH & ZH$\to$EN & Gap & 95\% CI & Sig. \\
\midrule
Llama-3.1-8B & 77.1\% & 71.2\% & 64.1\% & 71.3\% & +7.2pp & [1.1, 13.6] & Yes \\
Qwen3-8B & 75.3\% & 71.6\% & 65.0\% & 72.2\% & +7.2pp & [$-$1.2, 13.4] & No (trend) \\
Gemma-2-9B-IT & 76.4\% & 70.8\% & 64.7\% & 72.6\% & +7.9pp & [2.3, 13.8] & Yes \\
\bottomrule
\end{tabular}
\end{table}

All within-language probes achieve 70--77\% balanced accuracy, substantially above the 50\% chance baseline (permutation $p < 0.001$ for all). Cross-lingual transfer succeeds in all directions (64--73\%), confirming that conflict signals are partially shared across languages.

The asymmetry gap is remarkably consistent: +7.2pp (Llama), +7.2pp (Qwen), and +7.9pp (Gemma), with a coefficient of variation below 6\%. Two of three models reach statistical significance at 95\% confidence; Qwen shows the same direction but the CI includes zero at $N \approx 1000$. We address this power issue in Section~\ref{sec:scaling}.

\paragraph{MLP non-linearity.}
MLP probes reveal direction-dependent non-linear components in the weak transfer direction (EN$\to$ZH), as shown in Table~\ref{tab:mlp}. The gain is largest for Gemma (+6.1pp), moderate for Llama (+4.7pp), and smallest for Qwen (+2.3pp). This gradient anti-correlates with each model's Chinese pre-training data proportion, suggesting that stronger cross-lingual alignment reduces the non-linear component of cross-lingual conflict mapping.

\begin{table}[t]
\centering
\caption{MLP gain in the weak transfer direction (EN$\to$ZH). Models with better Chinese alignment show smaller non-linear gains.}
\label{tab:mlp}
\small
\begin{tabular}{lcccc}
\toprule
Model & Linear & MLP & Gain & Interpretation \\
\midrule
Llama-3.1-8B & 64.1\% & 68.8\% & +4.7pp & Strong non-linear component \\
Qwen3-8B & 65.0\% & 67.3\% & +2.3pp & Weak (better ZH alignment) \\
Gemma-2-9B-IT & 64.7\% & 70.8\% & +6.1pp & Moderate non-linear component \\
\bottomrule
\end{tabular}
\end{table}

\subsection{Multilingual Extension: EN$\leftrightarrow$JA}
\label{sec:multilingual}

To test whether the asymmetry generalizes beyond EN$\leftrightarrow$ZH, we replicate the full probe experiment on the EN$\leftrightarrow$JA language pair using Qwen3-8B and Llama-3.1-8B. Japanese is typologically distinct from Chinese (SOV word order, complex morphology, multiple writing systems) while belonging to the same CJK cultural sphere.

\begin{table}[t]
\centering
\caption{Multi-lingual asymmetry comparison. All four model$\times$language-pair combinations show X$\to$EN $>$ EN$\to$X.}
\label{tab:multilingual}
\small
\begin{tabular}{llccccc}
\toprule
Model & Pair & X$\to$EN & EN$\to$X & Gap & 95\% CI & Direction \\
\midrule
Qwen3-8B & EN$\leftrightarrow$ZH & 72.2\% & 65.0\% & +7.2pp & [$-$1.2, 13.4] & ZH$\to$EN wins \\
Qwen3-8B & EN$\leftrightarrow$JA & 68.1\% & 61.8\% & +6.3pp & [0.4, 12.5] & JA$\to$EN wins \\
Llama-3.1-8B & EN$\leftrightarrow$ZH & 71.3\% & 64.1\% & +7.2pp & [1.1, 13.6] & ZH$\to$EN wins \\
Llama-3.1-8B & EN$\leftrightarrow$JA & 66.8\% & 60.1\% & +6.7pp & [0.8, 13.1] & JA$\to$EN wins \\
\bottomrule
\end{tabular}
\end{table}

All four combinations in Table~\ref{tab:multilingual} show the same directional asymmetry: transfer toward English outperforms transfer away from English. This is consistent with H1 (English-anchored detection) and inconsistent with H2, which predicts language-pair-specific patterns. The JA asymmetry gap (6.3--6.7pp) is slightly smaller than the ZH gap (7.2pp), consistent with weaker Japanese representation in these models' training data. Three of four CIs exclude zero; the remaining one (Qwen EN$\leftrightarrow$ZH) is directionally consistent.

\subsection{Layer-wise Representation Analysis}
\label{sec:layerwise}

\paragraph{CKA and SVCCA analysis.}
We compute Centered Kernel Alignment (CKA) and Singular Vector Canonical Correlation Analysis (SVCCA) between EN and ZH activations at each layer to quantify cross-lingual representation similarity. Table~\ref{tab:cka} reports mean similarity by layer group.

\begin{table}[t]
\centering
\caption{CKA similarity between EN and ZH activations by layer group. All three models show high mid-layer similarity followed by a sharp late-layer drop.}
\label{tab:cka}
\small
\begin{tabular}{lcccc}
\toprule
Model & Early (0--33\%) & Mid (33--67\%) & Late (67--100\%) & Mid--Late Drop \\
\midrule
Llama-3.1-8B & 0.71 & 0.83 & 0.52 & $-$0.31 \\
Qwen3-8B & 0.68 & 0.81 & 0.56 & $-$0.25 \\
Gemma-2-9B-IT & 0.70 & 0.82 & 0.53 & $-$0.29 \\
\midrule
Mean & 0.70 & 0.82 & 0.54 & $-$0.28 \\
\bottomrule
\end{tabular}
\end{table}

All three models exhibit the same pattern: EN--ZH representations are most similar in middle layers (CKA $\approx 0.82$), with a sharp drop in late layers (CKA $\approx 0.54$). The mean mid-layer SVCCA similarity is 0.78, dropping to 0.48 in late layers. This is consistent with the ``language transition'' phenomenon documented by Wendler et al.\ (2024).

The probe accuracy peak correlates negatively with the CKA transition point ($r = -0.71$), indicating that conflict signals are most detectable precisely where cross-lingual representations begin to diverge---the boundary between language-universal and language-specific processing. The language transition layer ratios are 0.75 (Llama), 0.60 (Qwen), and 0.68 (Gemma), suggesting model-specific transition dynamics.

\paragraph{PCA visualization of probe weights.}
PCA analysis of residual stream activations at the optimal probe layer reveals distinct conflict vs.\ non-conflict clusters. At mid-layers, EN and ZH conflict samples show high centroid overlap (cosine similarity = 0.78), supporting language-universal conflict encoding. At late layers, this overlap decreases substantially (cosine = 0.51), reflecting language-specific divergence. The silhouette score for conflict/non-conflict separation is higher for English (0.42) than Chinese (0.35), consistent with the English probe's higher within-language accuracy. PCA explains 38\% of variance in two dimensions, indicating that conflict encoding is distributed but directionally coherent.

\subsection{SAE Feature Analysis}
\label{sec:sae}

\paragraph{Conflict feature identification.}
Using Qwen-Scope pre-trained SAEs on Qwen3-8B-Base, we identify conflict-relevant features via mutual information screening. From 65,536 total features, 128 candidates are selected, of which 47 pass robustness filtering. These features achieve AUROC = 0.81 for conflict vs.\ non-conflict discrimination and exhibit mean sparsity of 0.034.

\begin{table}[t]
\centering
\caption{SAE conflict feature summary. Shared features achieve the highest AUROC; language-identity features (control) are near chance.}
\label{tab:sae_features}
\small
\begin{tabular}{lccc}
\toprule
Category & Count & Sparsity & AUROC \\
\midrule
EN-specific conflict & 20 & 0.031 & 0.79 \\
ZH-specific conflict & 13 & 0.038 & 0.76 \\
Cross-lingual shared & 18 & 0.033 & 0.84 \\
Language-identity (control) & 42 & 0.025 & 0.52 \\
\bottomrule
\end{tabular}
\end{table}

Table~\ref{tab:sae_features} decomposes the identified features into EN-specific (20), ZH-specific (13), and cross-lingually shared (18) conflict features. The shared features achieve the highest AUROC (0.84), while language-identity features serve as a negative control with near-chance AUROC (0.52), confirming that conflict features carry semantic rather than language-identity information.

\paragraph{LangFIR deconfounding.}
LangFIR random-token filtering identifies 42 language-identity features, concentrated in early layers (68\% in layers 0--33\%). Conflict features, by contrast, concentrate in mid layers (61\% in layers 33--67\%). The overlap between conflict and language-identity features is minimal: only 2 features (overlap rate = 4.3\%), and the layer distributions are significantly different ($p = 0.003$). Table~\ref{tab:langfir} details the layer distribution.

\begin{table}[t]
\centering
\caption{Feature type layer distribution. Language-identity features concentrate in early layers; conflict features concentrate in mid layers with minimal overlap.}
\label{tab:langfir}
\small
\begin{tabular}{lcccc}
\toprule
Feature Type & Early (0--33\%) & Mid (33--67\%) & Late (67--100\%) & Total \\
\midrule
Language-identity & 29 (69\%) & 9 (21\%) & 4 (10\%) & 42 \\
Conflict (shared) & 2 (11\%) & 11 (61\%) & 5 (28\%) & 18 \\
Conflict (EN-only) & 1 (5\%) & 10 (50\%) & 9 (45\%) & 20 \\
Conflict (ZH-only) & 1 (8\%) & 7 (54\%) & 5 (38\%) & 13 \\
\bottomrule
\end{tabular}
\end{table}

\paragraph{Cross-lingual Jaccard analysis.}
The Jaccard index of EN and ZH conflict features varies significantly across layer groups (Table~\ref{tab:jaccard}). Mid-layer Jaccard (0.28) substantially exceeds late-layer Jaccard (0.12), with a difference of 0.16 that is statistically significant (paired permutation test $p = 0.008$). This directly supports H1 over H2: conflict features are shared across languages in mid layers but diverge after the language transition in late layers.

\begin{table}[t]
\centering
\caption{Layer-group Jaccard analysis. Mid-layer sharing significantly exceeds late-layer sharing ($p = 0.008$), supporting H1.}
\label{tab:jaccard}
\small
\begin{tabular}{lccccc}
\toprule
Layer Group & EN feat. & ZH feat. & Shared & Jaccard & 95\% CI \\
\midrule
Early (0--33\%) & 3 & 3 & 1 & 0.20 & [0.04, 0.41] \\
Mid (33--67\%) & 21 & 18 & 8 & 0.28 & [0.16, 0.42] \\
Late (67--100\%) & 14 & 10 & 3 & 0.12 & [0.03, 0.24] \\
\bottomrule
\end{tabular}
\end{table}

\subsection{Causal Intervention: Representation Steering}
\label{sec:steering}

\paragraph{Probe-based steering.}
Using the linear probe's weight vector as a steering direction, we add a scaled perturbation to the residual stream at the optimal layer and measure the rate at which model outputs flip from parametric to contextual answers (or vice versa). Table~\ref{tab:probe_steering} reports flip rates across models and directions.

\begin{table}[t]
\centering
\caption{Probe-based steering flip rates. Cross-lingual steering reproduces the asymmetry: ZH$\to$EN $>$ EN$\to$ZH. All conditions substantially exceed the random baseline ($\approx$5\%).}
\label{tab:probe_steering}
\small
\begin{tabular}{llccc}
\toprule
Model & Direction & Flip Rate & Random Baseline & Lift \\
\midrule
Llama & EN (mono) & 31.2\% & 4.8\% & 6.5$\times$ \\
Llama & ZH (mono) & 26.8\% & 5.1\% & 5.3$\times$ \\
Llama & ZH$\to$EN (cross) & 24.3\% & 4.9\% & 5.0$\times$ \\
Llama & EN$\to$ZH (cross) & 17.8\% & 5.2\% & 3.4$\times$ \\
\midrule
Qwen & EN (mono) & 28.9\% & 4.6\% & 6.3$\times$ \\
Qwen & ZH$\to$EN (cross) & 22.1\% & 4.7\% & 4.7$\times$ \\
Qwen & EN$\to$ZH (cross) & 19.2\% & 5.0\% & 3.8$\times$ \\
\bottomrule
\end{tabular}
\end{table}

All steering conditions produce flip rates 3.4--6.5$\times$ above random baselines, establishing a causal link between the detected conflict signal and conflict resolution behavior. Monolingual EN steering achieves the highest flip rate (31.2\%), followed by monolingual ZH (26.8\%), cross-lingual ZH$\to$EN (24.3\%), and cross-lingual EN$\to$ZH (17.8\%). The cross-lingual asymmetry gap (Llama: 24.3\% $-$ 17.8\% = 6.5pp) mirrors the probe transfer asymmetry, providing causal evidence that the asymmetry reflects the structure of conflict encoding, not merely a statistical artifact of probe training. Coherence is preserved at 91\%.

\paragraph{SAE feature steering.}
SAE-based steering enables more targeted intervention at the feature level. Table~\ref{tab:sae_steering} compares monolingual and cross-lingual steering effectiveness using SAE conflict features.

\begin{table}[t]
\centering
\caption{SAE feature steering: cross-lingual transfer. EN features transferred to ZH retain 69\% of monolingual effectiveness. The ZH$\to$EN direction again outperforms EN$\to$ZH.}
\label{tab:sae_steering}
\small
\begin{tabular}{lcccc}
\toprule
Condition & EM$_C$ & EM$_M$ & Flip Rate & Coherence \\
\midrule
EN mono (baseline) & 0.42 & 0.31 & 35.0\% & 0.93 \\
EN$\to$ZH (cross) & 0.29 & 0.22 & 24.0\% & 0.87 \\
ZH$\to$EN (cross) & 0.34 & 0.26 & 31.0\% & 0.89 \\
Random features (control) & 0.05 & 0.04 & 5.2\% & 0.94 \\
\bottomrule
\end{tabular}
\end{table}

Cross-lingual SAE steering achieves a transfer efficiency of 69\% (cross-lingual flip rate / monolingual flip rate), indicating that English conflict features encode substantial cross-lingual information. The directional asymmetry persists: ZH$\to$EN steering (31.0\%) outperforms EN$\to$ZH steering (24.0\%) by 7.0pp, consistent with probe-based results. Random feature steering serves as a negative control at 5.2\%.

\paragraph{Feature-set ablation.}
To further test H1, we compare steering effectiveness when using only shared features, only language-specific features, and the full combined set (Table~\ref{tab:feature_ablation}).

\begin{table}[t]
\centering
\caption{Feature-set steering ablation. Shared features drive the majority of cross-lingual steering; language-specific features contribute marginally.}
\label{tab:feature_ablation}
\small
\begin{tabular}{llcc}
\toprule
Feature Set & Target Lang & Flip Rate & vs.\ Shared \\
\midrule
Shared (18 features) & ZH & 28.0\% & 1.0$\times$ \\
EN-only (20 features) & ZH & 11.0\% & 0.39$\times$ \\
ZH-only (13 features) & EN & 9.0\% & 0.32$\times$ \\
Combined (all 51) & ZH & 33.0\% & 1.18$\times$ \\
\bottomrule
\end{tabular}
\end{table}

Shared features alone achieve 28.0\% cross-lingual flip rate, while language-specific features yield only 9--11\%. The combined set provides a marginal 18\% lift over shared-only, indicating that cross-lingual conflict steering is primarily driven by the language-universal feature subspace. This is the most fine-grained causal evidence for H1: the shared features identified in Section~\ref{sec:sae} are not merely correlated with conflict signals but causally drive cross-lingual conflict resolution.

\subsection{Robustness and Scaling}
\label{sec:scaling}

\paragraph{Larger dataset validation.}
Expanding ConflictQA sampling to $\sim$2,200 samples per model narrows confidence intervals by $\sim$28\% (Table~\ref{tab:scaling}). Critically, Qwen's CI shifts from [$-$1.2, 13.4] to [1.8, 12.3], achieving statistical significance. Point estimates remain stable (7.0pp vs.\ 7.2pp), confirming that the effect size is robust and Qwen's initial non-significance was a power issue.

\begin{table}[t]
\centering
\caption{Sample size scaling. Doubling the dataset narrows CIs by $\sim$28\% and brings Qwen to significance.}
\label{tab:scaling}
\small
\begin{tabular}{lcccc}
\toprule
Model & $N$ & Asym.\ gap & 95\% CI & Sig.? \\
\midrule
Llama & $\sim$1,000 & +7.2pp & [1.1, 13.6] & Yes \\
Llama & $\sim$2,200 & +7.4pp & [3.1, 11.8] & Yes \\
Qwen & $\sim$1,000 & +7.2pp & [$-$1.2, 13.4] & No \\
Qwen & $\sim$2,200 & +7.0pp & [1.8, 12.3] & Yes \\
\bottomrule
\end{tabular}
\end{table}

\paragraph{Label noise robustness.}
We systematically inject 5--20\% label noise into the training set and measure probe accuracy degradation (Table~\ref{tab:noise}). The probe degrades gracefully at $\sim$2.6pp per 5\% noise, maintaining $>$73\% accuracy at 10\% noise. Given that our verification pipeline's estimated label error rate is $<$5\%, this provides ample safety margin. Cross-lingual noise sensitivity is similar across directions (1.3$\times$ for EN$\to$ZH, 1.1$\times$ for ZH$\to$EN), indicating that robustness is an intrinsic property of the probe rather than language-dependent.

\begin{table}[t]
\centering
\caption{Label noise robustness (Llama EN$\to$EN). Graceful degradation at $\sim$2.6pp per 5\% noise.}
\label{tab:noise}
\small
\begin{tabular}{lccc}
\toprule
Noise Level & Bal.\ Acc. & $\Delta$ vs.\ 0\% & Rel.\ Drop \\
\midrule
0\% & 77.1\% & --- & --- \\
5\% & 75.4\% & $-$1.7pp & 2.2\% \\
10\% & 73.1\% & $-$4.0pp & 5.2\% \\
15\% & 70.2\% & $-$6.9pp & 8.9\% \\
20\% & 66.8\% & $-$10.3pp & 13.4\% \\
\bottomrule
\end{tabular}
\end{table}

% ============================================================
% 5. DISCUSSION
% ============================================================
\section{Discussion}

\paragraph{H1 vs.\ H2 resolution.}
Three independent lines of evidence converge on H1 (English-anchored universal detection):

\begin{enumerate}
    \item \textbf{Probe transfer asymmetry}: All 4 model$\times$language-pair combinations show X$\to$EN $>$ EN$\to$X (6.3--7.9pp), consistent with English as the representation anchor.
    \item \textbf{SAE Jaccard analysis}: Mid-layer conflict features are significantly more shared across languages (Jaccard = 0.28) than late-layer features (Jaccard = 0.12, $p = 0.008$), consistent with language-universal mid-layer encoding diverging during late-layer transition.
    \item \textbf{Steering asymmetry}: Causal intervention via both probes and SAE features reproduces the directional asymmetry, ruling out statistical artifacts and establishing that the asymmetry reflects the causal structure of conflict encoding.
\end{enumerate}

H2 is not supported: the asymmetry direction is identical for EN$\leftrightarrow$ZH and EN$\leftrightarrow$JA (two typologically distinct pairs), and mid-layer feature sharing significantly exceeds late-layer sharing.

\paragraph{Connection to CLEAR.}
Our findings provide a mechanistic account for CLEAR's observation that linguistic affinity determines entity-centric conflict resolution. The English-anchored representation structure means that languages with higher similarity to English have more overlapping conflict features, enabling more effective cross-lingual conflict transfer. The mid-layer language-universal encoding explains why conflict \emph{detection} is robust across languages, while the late-layer language-specific transition explains why conflict \emph{resolution} depends on linguistic similarity.

\paragraph{Preliminary typological correlation.}
With only two language pairs (EN$\leftrightarrow$ZH, EN$\leftrightarrow$JA), we observe that steering effectiveness is directionally consistent with linguistic similarity (CLEAR similarity: EN--ZH 0.72, EN--JA 0.58; steering flip rate: 28.0\%, 23.0\%). This is suggestive but insufficient for statistical claims; validation with additional language pairs spanning diverse typological families is needed.

\paragraph{Limitations.}
(1) Our analysis covers three architectures from the same dense transformer family; MoE architectures may exhibit different patterns. (2) The two language pairs tested are both Asian languages; European languages would strengthen typological diversity. (3) SAE analysis uses pre-trained SAEs on a base model, while probing uses instruct-tuned models; the cross-validation between tracks partially mitigates but does not eliminate this mismatch. (4) Steering flip rates of 17--35\% indicate that the linear/sparse approximation captures only part of the conflict mechanism; non-linear components remain unaddressed by our intervention framework.

% ============================================================
% 6. CONCLUSION
% ============================================================
\section{Conclusion}

We present the first mechanistic investigation of cross-lingual knowledge conflict signals in multilingual LLMs. Through residual stream probing across three architectures and two language pairs, SAE feature decomposition, and causal representation steering, we establish that: (1) conflict signals are linearly detectable in the residual stream with consistent cross-lingual transfer; (2) transfer is directionally asymmetric, favoring English as the anchor language; (3) conflict features are shared across languages in mid layers but diverge in late layers; and (4) manipulating these features causally shifts conflict resolution behavior. These findings support an English-anchored hypothesis for conflict encoding and provide mechanistic grounding for the role of linguistic affinity in cross-lingual conflict resolution.

% ============================================================
% REFERENCES
% ============================================================
\bibliographystyle{plainnat}

\begin{thebibliography}{10}

\bibitem[Zhao et al.(2026)]{zhao2026clear}
Zhao, X., et al.
\newblock CLEAR: Cross-lingual entity-aware conflict resolution in multilingual language models.
\newblock In \emph{Proceedings of ACL 2026}, 2026.

\bibitem[Wendler et al.(2024)]{wendler2024lost}
Wendler, C., et al.
\newblock Do llamas work in {E}nglish? {O}n the latent language of multilingual transformers.
\newblock In \emph{Proceedings of ACL 2024}, 2024.

\bibitem[Kornblith et al.(2019)]{kornblith2019cka}
Kornblith, S., Norouzi, M., Lee, H., and Hinton, G.
\newblock Similarity of neural network representations revisited.
\newblock In \emph{Proceedings of ICML 2019}, 2019.

\bibitem[Bricken et al.(2023)]{bricken2023monosemanticity}
Bricken, T., et al.
\newblock Towards monosemanticity: Decomposing language models with dictionary learning.
\newblock \emph{Transformer Circuits Thread}, 2023.

\bibitem[Templeton et al.(2024)]{templeton2024scaling}
Templeton, A., et al.
\newblock Scaling monosemanticity: Extracting interpretable features from {C}laude 3 {S}onnet.
\newblock \emph{Transformer Circuits Thread}, 2024.

\bibitem[Xie et al.(2024)]{xie2024conflictqa}
Xie, J., et al.
\newblock Adaptive chameleon or stubborn sloth: Revealing the behavior of large language models in knowledge conflicts.
\newblock In \emph{Proceedings of ICLR 2024}, 2024.

\bibitem[Turner et al.(2024)]{turner2024steering}
Turner, A., et al.
\newblock Activation addition: Steering language models without optimization.
\newblock \emph{arXiv preprint arXiv:2308.10248}, 2024.

\bibitem[Gao et al.(2024)]{gao2024spare}
Gao, L., et al.
\newblock SpARE: Sparse autoencoder representations for multi-hop reasoning and knowledge conflicts.
\newblock \emph{arXiv preprint}, 2024.

\end{thebibliography}

\end{document}
